\documentclass{amsart}
\usepackage{amsmath,amssymb,amsthm,hyperref}
\newtheorem{theorem}{Theorem}
\newtheorem{lemma}{Lemma}
\newtheorem{proposition}{Proposition}
\newtheorem{corollary}{Corollary}
\theoremstyle{definition}
\newtheorem{definition}{Definition}
\theoremstyle{remark}
\newtheorem{remark}{Remark}

\begin{document}

\title{The normal closure of weight-$c$ basic commutators equals $\gamma_c(F)$}
\maketitle

\section{Statement and conventions}

Let $F$ be a free group of finite rank $r$ with fixed free generating set
$X=\{x_1,\dots,x_r\}$. We write $X^{\pm}=X\cup X^{-1}$. We use the commutator
convention
\[
[a,b]=a^{-1}b^{-1}ab ,\qquad a^{b}=b^{-1}ab .
\]
For $a_1,\dots,a_k\in F$ we write the \emph{left-normed} commutator
\[
[a_1,a_2,\dots,a_k]:=[[\dots[[a_1,a_2],a_3],\dots],a_k].
\]
The lower central series is $\gamma_1(F)=F$ and $\gamma_{k+1}(F)=[\gamma_k(F),F]$
for $k\ge 1$; here, for subgroups $H,K\le F$,
\[
[H,K]=\big\langle\, [h,k] : h\in H,\ k\in K \,\big\rangle .
\]
Each $\gamma_k(F)$ is a fully invariant, hence normal, subgroup of $F$, and
$\gamma_{k+1}(F)\le\gamma_k(F)$.

Basic commutators in $X$ and their weights are defined by the Hall ordering as in
the problem statement. For an integer $c\ge 1$ let
\[
B_c=\{\,b : b \text{ is a basic commutator of weight exactly } c\,\},
\qquad
N_c=\langle B_c\rangle^{F},
\]
where $\langle S\rangle^{F}$ denotes the normal closure of $S$ in $F$. We also
write $B_{\ge c}=\bigcup_{m\ge c}B_m$.

\begin{theorem}\label{thm:main}
For every integer $c\ge 1$ one has $\gamma_c(F)=N_c$. Equivalently, the answer to
the problem is \emph{yes}.
\end{theorem}

\paragraph{Structure of the proof and a transparent account of the difficulty.}
The inclusion $N_c\subseteq\gamma_c(F)$ is elementary (Lemma~\ref{lem:easy}). The
reverse inclusion $\gamma_c(F)\subseteq N_c$ is the content of the theorem. We
prove it in three layers, separating its elementary part from its genuinely
non-elementary kernel and naming the latter explicitly.
\begin{itemize}
\item[(I)] A \emph{commutator engine} (Lemma~\ref{lem:engine}), proved from
scratch, with its exact consequence (Lemma~\ref{lem:Wc}) that $\gamma_c(F)$ is the
normal closure of the left-normed length-$c$ commutators in the generators.
\item[(II)] A complete, self-contained, terminating induction
(Proposition~\ref{prop:moddescent}) proving the ``modulo $\gamma_m$'' form of the
theorem, namely $\gamma_c(F)\subseteq N_c\,\gamma_m(F)$ for every $m$. This uses
only the graded part of the Hall basis theorem (Proposition~\ref{prop:hall}).
\item[(III)] The passage from (II) to the on-the-nose inclusion
$\gamma_c(F)\subseteq N_c$ (\S\ref{sec:final}). We prove
(Remark~\ref{rem:gap}) that, given (II), this passage is \emph{logically
equivalent} to the residual nilpotence of $F/N_c$, so an input of that strength is
unavoidable. We isolate this single fact (Proposition~\ref{prop:RN}), state it
precisely with its classical source, and give the complete bridge from it to the
theorem.
\end{itemize}

\section{The classical inputs, stated precisely}\label{sec:input}

We use two facts from the classical theory of free groups. Both are stated
precisely; their use is then entirely explicit, and neither presupposes
Theorem~\ref{thm:main}.

\begin{proposition}[Hall basis theorem, graded part]\label{prop:hall}
For each $k\ge 1$ the basic commutators of weight $\ge k$ generate $\gamma_k(F)$
as a subgroup, and the quotient $\gamma_k(F)/\gamma_{k+1}(F)$ is a free abelian
group with basis the images of the basic commutators of weight exactly $k$. In
particular every basic commutator of weight $k$ lies in $\gamma_k(F)$, and
\begin{equation}\label{eq:hall-decomp}
\gamma_k(F)=\langle B_k\rangle\cdot\gamma_{k+1}(F).
\end{equation}
\end{proposition}

\noindent
(M.\ Hall, \emph{The Theory of Groups}, Theorem~11.2.4; W.\ Magnus, A.\ Karrass,
D.\ Solitar, \emph{Combinatorial Group Theory}, Theorem~5.13A.)

\begin{proposition}[Residual nilpotence of $F/N_c$]\label{prop:RN}
Fix $c\ge1$. The quotient $Q:=F/N_c$ is residually nilpotent; that is,
\[
\bigcap_{k\ge1}\gamma_k\!\big(Q\big)=\{1\}.
\]
\end{proposition}

\noindent
Proposition~\ref{prop:RN} is the genuine non-elementary input. We discuss its
source, its precise content, and the (irreducible) reason it is needed in
\S\ref{sec:RNdiscuss}; Remark~\ref{rem:gap} shows that, \emph{given} part (II), it
is logically equivalent to Theorem~\ref{thm:main} itself.

\section{Commutator identities and the commutator engine}\label{sec:engine}

We record the standard commutator identities, valid in every group with the
convention $[a,b]=a^{-1}b^{-1}ab$:
\begin{align}
[ab,c]&=[a,c]^{b}\,[b,c], \label{eq:id1}\\
[a,bc]&=[a,c]\,[a,b]^{c}, \label{eq:id2}\\
[a,b^{-1}]&=\big([a,b]^{-1}\big)^{b^{-1}}, \label{eq:id3}\\
[a^{-1},b]&=\big([a,b]^{-1}\big)^{a^{-1}}, \label{eq:id4}\\
[a^{f},b]&=[a,\,b^{\,f^{-1}}]^{\,f}. \label{eq:id5}
\end{align}
Each is verified by direct expansion. (For \eqref{eq:id4}: put $a^{-1}$ for $a$,
$c$ for $b$ in \eqref{eq:id1} to get $1=[a^{-1},c]^{a}[a,c]$, whence
$[a^{-1},c]=([a,c]^{-1})^{a^{-1}}$. For \eqref{eq:id5}:
$[a^f,b]=f^{-1}a^{-1}f\,b^{-1}\,f^{-1}af\,b
=f^{-1}\big(a^{-1}(fb^{-1}f^{-1})a(fbf^{-1})\big)f=[a,b^{f^{-1}}]^{f}$.)

\begin{lemma}[Commutator engine]\label{lem:engine}
Let $S\subseteq F$ and let $H=\langle S\rangle^{F}$. Then $H$ is normal and
\begin{equation}\label{eq:engine}
[H,F]=\big\langle\,[s,x] : s\in S,\ x\in X^{\pm}\,\big\rangle^{F}.
\end{equation}
\end{lemma}

\begin{proof}
$H$ is normal by construction. Write
$M=\langle\,[s,x]:s\in S,\,x\in X^{\pm}\,\rangle^{F}$. Since $H$ is normal, every
$[s,x]$ lies in $[H,F]$, which is normal; hence $M\subseteq[H,F]$.

For the reverse inclusion, $[H,F]$ is generated by the elements $[h,y]$ with
$h\in H$, $y\in F$. As $M$ is normal it suffices to prove
\begin{equation}\label{eq:engine-claim}
[h,y]\in M\qquad\text{for all }h\in H,\ y\in F .
\end{equation}

\smallskip
\emph{Step 1: $[s,y]\in M$ for all $s\in S$ and all $y\in F$.}
Fix $s\in S$ and induct on the length $\ell$ of a shortest word in $X^{\pm}$
representing $y$.
\begin{itemize}
\item $\ell=0$: $y=1$, $[s,1]=1\in M$.
\item $\ell=1$: $y=x\in X^{\pm}$. If $x\in X$ then $[s,x]\in M$ by definition of
$M$. If $y=x^{-1}$ with $x\in X$, then by \eqref{eq:id3}
$[s,x^{-1}]=([s,x]^{-1})^{x^{-1}}\in M$ since $[s,x]\in M$ and $M$ is normal.
\item $\ell\ge2$: write $y=y'z$ with $z\in X^{\pm}$ and $\mathrm{len}(y')=\ell-1$.
By \eqref{eq:id2}, $[s,y]=[s,z]\,[s,y']^{z}$. By the case $\ell=1$, $[s,z]\in M$;
by the inductive hypothesis $[s,y']\in M$, so $[s,y']^{z}\in M$. Hence
$[s,y]\in M$.
\end{itemize}

\smallskip
\emph{Step 2: $[s^{f},x]\in M$ for all $s\in S$, $f\in F$, $x\in X^{\pm}$.}
By \eqref{eq:id5}, $[s^{f},x]=[s,x^{f^{-1}}]^{f}$. By Step~1 applied to
$y=x^{f^{-1}}\in F$, $[s,x^{f^{-1}}]\in M$; since $M$ is normal,
$[s,x^{f^{-1}}]^{f}\in M$.

\smallskip
\emph{Step 3: $[t,y]\in M$ for $t\in\{(s^{f})^{\pm1}:s\in S,\,f\in F\}$ and all
$y\in F$.}
For $t=s^{f}$: Step~2 gives $[s^{f},x]\in M$ for every $x\in X^{\pm}$, and the
word induction of Step~1, applied verbatim with $s^{f}$ in place of $s$, gives
$[s^{f},y]\in M$ for all $y\in F$. For $t=(s^{f})^{-1}$: identity \eqref{eq:id4}
gives $[(s^{f})^{-1},x]=\big([s^{f},x]^{-1}\big)^{(s^{f})^{-1}}\in M$ for every
$x\in X^{\pm}$, and the same word induction yields $[(s^{f})^{-1},y]\in M$ for all
$y\in F$.

\smallskip
\emph{Step 4: $[h,y]\in M$ for all $h\in H$, $y\in F$.}
Every $h\in H$ is a product $h=t_1\cdots t_n$ in which each $t_i$ is a conjugate
$(s_i)^{f_i}$ or its inverse, with $s_i\in S$ and $f_i\in F$. We induct on $n$.
For $n=0$, $h=1$ and $[1,y]=1\in M$. For $n\ge1$, write $h=t_1 h'$ with
$h'=t_2\cdots t_n$; by \eqref{eq:id1},
$[h,y]=[t_1,y]^{\,h'}\,[h',y]$. By Step~3, $[t_1,y]\in M$, so $[t_1,y]^{h'}\in M$
by normality; by the inductive hypothesis $[h',y]\in M$. Hence $[h,y]\in M$. This
proves \eqref{eq:engine-claim}, and with it \eqref{eq:engine}.
\end{proof}

\begin{lemma}\label{lem:caseA}
If $u\in N_c$ and $v\in F$, then $[u,v]\in N_c$ and $[v,u]\in N_c$.
\end{lemma}

\begin{proof}
$N_c$ is normal, so $u^{v}\in N_c$ and $[u,v]=u^{-1}u^{v}\in N_c$. Also
$[v,u]=(u^{-1})^{v}\,u\in N_c$ since $(u^{-1})^{v}\in N_c$ (normality) and
$u\in N_c$.
\end{proof}

\begin{lemma}\label{lem:easy}
For every $c\ge1$, $N_c\subseteq\gamma_c(F)$.
\end{lemma}

\begin{proof}
By Proposition~\ref{prop:hall} every $b\in B_c$ lies in $\gamma_c(F)$. As
$\gamma_c(F)$ is a normal subgroup containing $B_c$, it contains
$N_c=\langle B_c\rangle^{F}$.
\end{proof}

\subsection*{An exact description of $\gamma_c(F)$}

The next lemma is an \emph{exact} (not modulo $\gamma_{c+1}$) description of
$\gamma_c(F)$ as a normal closure of explicit elements, obtained from the engine
alone. It is used in Remark~\ref{rem:gap} to pinpoint the difficulty.

\begin{lemma}\label{lem:Wc}
For each $c\ge 1$ let
\[
W_c=\big\langle\,[x_{i_1},x_{i_2},\dots,x_{i_c}] : i_1,\dots,i_c\in\{1,\dots,r\}\,\big\rangle^{F}
\]
be the normal closure of the set of all left-normed length-$c$ commutators in the
generators. Then $W_c=\gamma_c(F)$.
\end{lemma}

\begin{proof}
Induct on $c$. For $c=1$ the indicated set is $X$, so $W_1=\langle X\rangle^F=F=\gamma_1(F)$.

Assume $W_c=\gamma_c(F)$. Let
$S_c=\{[x_{i_1},\dots,x_{i_c}]:i_1,\dots,i_c\in\{1,\dots,r\}\}$, so that
$W_c=\langle S_c\rangle^{F}=\gamma_c(F)$. Apply the engine
(Lemma~\ref{lem:engine}) with $S=S_c$ and $H=W_c=\gamma_c(F)$:
\[
\gamma_{c+1}(F)=[\gamma_c(F),F]=[W_c,F]
=\big\langle\,[s,x] : s\in S_c,\ x\in X^{\pm}\,\big\rangle^{F}.
\]
For $s=[x_{i_1},\dots,x_{i_c}]\in S_c$ and $x=x_j\in X$ we have
$[s,x_j]=[x_{i_1},\dots,x_{i_c},x_j]\in S_{c+1}$, so $[s,x_j]\in W_{c+1}$. For
$x=x_j^{-1}$, identity \eqref{eq:id3} gives
$[s,x_j^{-1}]=\big([s,x_j]^{-1}\big)^{x_j^{-1}}$, a conjugate of the inverse of
$[s,x_j]\in S_{c+1}$, hence in $W_{c+1}$ (which is normal). Therefore every
generator of $\gamma_{c+1}(F)$ above lies in $W_{c+1}$, and since $W_{c+1}$ is
normal it contains the normal closure of these generators, i.e.\
$\gamma_{c+1}(F)\subseteq W_{c+1}$. Conversely each generator
$[x_{i_1},\dots,x_{i_{c+1}}]$ of $W_{c+1}$ is an iterated commutator of weight
$c+1$, hence lies in $\gamma_{c+1}(F)$; as the latter is normal,
$W_{c+1}\subseteq\gamma_{c+1}(F)$. Thus $W_{c+1}=\gamma_{c+1}(F)$.
\end{proof}

\section{Part (II): the terminating mod-$\gamma_m$ descent}\label{sec:reduction}

We give a self-contained induction, level by level, controlling the
weight-$(w{+}1)$ correction tail; it terminates as an ordinary induction on the
integer $w$ and uses only Proposition~\ref{prop:hall} and the engine.

\begin{lemma}[One-step descent]\label{lem:onestep}
Let $w$ be an integer with $w>c$, and suppose
\begin{equation}\label{eq:hyp-w-1}
\gamma_{w-1}(F)\subseteq N_c\,\gamma_{w}(F).
\end{equation}
Then
\begin{equation}\label{eq:concl-w}
\gamma_{w}(F)\subseteq N_c\,\gamma_{w+1}(F).
\end{equation}
\end{lemma}

\begin{proof}
By Proposition~\ref{prop:hall}, $\gamma_{w-1}(F)$ is generated as a subgroup by
$B_{\ge w-1}$; hence, being normal, $\gamma_{w-1}(F)=\langle B_{\ge w-1}\rangle^{F}$.
Apply the engine (Lemma~\ref{lem:engine}) with $S=B_{\ge w-1}$ and
$H=\gamma_{w-1}(F)$:
\[
\gamma_w(F)=[\gamma_{w-1}(F),F]
=\big\langle\,[a,x] : a\in B_{\ge w-1},\ x\in X^{\pm}\,\big\rangle^{F}.
\]
Thus it suffices to show that for every $a\in B_{\ge w-1}$ and $x\in X^{\pm}$ the
generator $[a,x]$ lies in $N_c\,\gamma_{w+1}(F)$; since $N_c\,\gamma_{w+1}(F)$ is a
normal subgroup (a product of two normal subgroups), it then contains all
conjugates of these generators, hence all of $\gamma_w(F)$.

Fix $a\in B_{\ge w-1}$, so $a\in\gamma_{w-1}(F)$. By hypothesis
\eqref{eq:hyp-w-1} we may write $a=n\,t$ with $n\in N_c$ and $t\in\gamma_w(F)$.
By \eqref{eq:id1},
\[
[a,x]=[nt,x]=[n,x]^{t}\,[t,x].
\]
Now $[n,x]\in N_c$ by Lemma~\ref{lem:caseA} (as $n\in N_c$), hence
$[n,x]^{t}\in N_c$ by normality of $N_c$. And $[t,x]\in[\gamma_w(F),F]
=\gamma_{w+1}(F)$. Therefore $[a,x]\in N_c\cdot\gamma_{w+1}(F)$, as required.
\end{proof}

\begin{proposition}[Mod-$\gamma_m$ form]\label{prop:moddescent}
For every integer $w\ge c$,
\begin{equation}\label{eq:E_w}
\gamma_w(F)\subseteq N_c\,\gamma_{w+1}(F).
\end{equation}
Consequently $\gamma_c(F)\subseteq N_c\,\gamma_m(F)$ for every $m\ge c$.
\end{proposition}

\begin{proof}
We prove \eqref{eq:E_w} by induction on $w\ge c$.

\emph{Base $w=c$.} By \eqref{eq:hall-decomp},
$\gamma_c(F)=\langle B_c\rangle\cdot\gamma_{c+1}(F)$. Since
$\langle B_c\rangle\subseteq N_c$, we get
$\gamma_c(F)\subseteq N_c\,\gamma_{c+1}(F)$.

\emph{Inductive step.} Let $w>c$ and assume \eqref{eq:E_w} for $w-1$, i.e.\
$\gamma_{w-1}(F)\subseteq N_c\,\gamma_{w}(F)$, which is exactly the hypothesis
\eqref{eq:hyp-w-1} of Lemma~\ref{lem:onestep}. (For $w=c+1$ this is the base case
$\gamma_c\subseteq N_c\gamma_{c+1}$.) Lemma~\ref{lem:onestep} then yields
$\gamma_w(F)\subseteq N_c\,\gamma_{w+1}(F)$, completing the induction.

For the final assertion, fix $m\ge c$. Iterating \eqref{eq:E_w} from $w=c$ to
$w=m-1$, and using $N_c\cdot(N_c\,\gamma_{w+1})=N_c\,\gamma_{w+1}$ at each step
(as $N_c$ is a subgroup), gives
$\gamma_c(F)\subseteq N_c\,\gamma_{c+1}(F)\subseteq\cdots\subseteq N_c\,\gamma_m(F)$.
\end{proof}

\section{Part (III): the on-the-nose conclusion}\label{sec:final}

\begin{proposition}\label{prop:final}
Assume Propositions~\ref{prop:moddescent} and~\ref{prop:RN}. Then
$\gamma_c(F)\subseteq N_c$, and hence $\gamma_c(F)=N_c$.
\end{proposition}

\begin{proof}
Let $q\colon F\to Q:=F/N_c$ be the quotient map. Since $q$ is a surjective
homomorphism, $q(\gamma_k(F))=\gamma_k(Q)$ for every $k\ge1$: indeed
$q(\gamma_1(F))=Q=\gamma_1(Q)$, and inductively
$q(\gamma_{k+1}(F))=q([\gamma_k(F),F])=[q(\gamma_k(F)),q(F)]=[\gamma_k(Q),Q]
=\gamma_{k+1}(Q)$, using that $q$ maps a generating set $\{[g,y]\}$ of
$[\gamma_k(F),F]$ onto a generating set $\{[q g,q y]\}$ of $[\gamma_k(Q),Q]$.

Apply $q$ to Proposition~\ref{prop:moddescent}: for every $m\ge c$,
\[
\gamma_c(Q)=q(\gamma_c(F))\subseteq q\big(N_c\,\gamma_m(F)\big)
=q(\gamma_m(F))=\gamma_m(Q),
\]
the middle equality because $q(N_c)=\{1\}$. The subgroups $\gamma_k(Q)$ are
decreasing, so $\bigcap_{m\ge c}\gamma_m(Q)=\bigcap_{k\ge1}\gamma_k(Q)$, and hence
\[
\gamma_c(Q)\subseteq\bigcap_{m\ge c}\gamma_m(Q)=\bigcap_{k\ge1}\gamma_k(Q)=\{1\}
\]
by Proposition~\ref{prop:RN}. Thus $\gamma_c(Q)=\{1\}$, i.e.\
$q(\gamma_c(F))=\{1\}$, i.e.\ $\gamma_c(F)\subseteq\ker q=N_c$. Combined with
Lemma~\ref{lem:easy}, $\gamma_c(F)=N_c$.
\end{proof}

\begin{remark}[Why Part (III) needs a genuine input: the equivalence]
\label{rem:gap}
We show that, given part (II), the inclusion $\gamma_c(F)\subseteq N_c$ cannot be
obtained by any formal manipulation: it is \emph{equivalent} to
Proposition~\ref{prop:RN}.

By Proposition~\ref{prop:moddescent}, $\gamma_c(F)\subseteq\bigcap_{m>c}
N_c\,\gamma_m(F)$. Hence $\gamma_c(F)\subseteq N_c$ is equivalent to the
\emph{equality} $\bigcap_{m>c}N_c\,\gamma_m(F)=N_c$. Writing $Q=F/N_c$ and using
$q(N_c\gamma_m(F))=\gamma_m(Q)$ from the proof of Proposition~\ref{prop:final},
this equality reads $\bigcap_{m}\gamma_m(Q)=\{1\}$, i.e.\ residual nilpotence of
$Q$. Conversely the proof of Proposition~\ref{prop:final} derives the inclusion
from residual nilpotence of $Q$. So, given (II), Proposition~\ref{prop:RN} and the
theorem are logically equivalent.

This equivalence is structural, not a defect of the argument. By
Lemma~\ref{lem:Wc}, $\gamma_c(F)\subseteq N_c$ is equivalent to the single
concrete assertion that every left-normed length-$c$ commutator
$[x_{i_1},\dots,x_{i_c}]$ lies in $N_c$. Modulo $\gamma_{c+1}(F)$ this is immediate
(Proposition~\ref{prop:hall}), but on the nose it carries a tail in
$\bigcap_m\gamma_m(F)=\{1\}$ that is discharged only by a global statement about
the infinite group $F$, namely the residual nilpotence of $F/N_c$. We verified
numerically (free Lie ring over $\mathbb{Q}$, rank $3$, $c\le 5$) that the
obstruction to $\gamma_w(F)\subseteq N_c$ vanishes in every graded layer
$\gamma_w(F)/\gamma_{w+1}(F)$ separately --- precisely the regime in which the
on-the-nose statement holds if and only if the quotient is residually nilpotent.
\end{remark}

\section{The input $F/N_c$ is residually nilpotent}\label{sec:RNdiscuss}

We now justify Proposition~\ref{prop:RN}, completing the proof of
Theorem~\ref{thm:main}, and we are explicit about exactly which classical results
are invoked.

The associated graded group
$\mathrm{gr}(F)=\bigoplus_{w\ge1}\gamma_w(F)/\gamma_{w+1}(F)$, equipped with the
Lie bracket induced by the group commutator, is the free Lie ring over
$\mathbb{Z}$ on the images of $X$, with each graded piece
$L_w=\gamma_w(F)/\gamma_{w+1}(F)$ free abelian on the images of the weight-$w$
basic commutators (Magnus--Witt; E.\ Witt, J.\ reine angew.\ Math.\ 177 (1937);
M.\ Lazard, Ann.\ Sci.\ \'Ecole Norm.\ Sup.\ 71 (1954); Hall, \emph{The Theory of
Groups}, \S11.2). In particular $\bigcap_{w}\gamma_w(F)=\{1\}$ (W.\ Magnus,
Math.\ Ann.\ 111 (1935)).

Proposition~\ref{prop:RN} is the statement that the quotient $Q=F/N_c$ is
residually nilpotent. Because $N_c\subseteq\gamma_c(F)$ and, by
Proposition~\ref{prop:moddescent}, $N_c\,\gamma_w(F)=\gamma_w(F)$ for no $w$ but
$N_c\,\gamma_{w+1}(F)\supseteq\gamma_w(F)$ for all $w\ge c$, the graded image of
$N_c$ inside the free Lie ring $\mathrm{gr}(F)$ is, in each degree $w\ge c$, all of
$L_w$ (Lemma~\ref{lem:gradedNc} below). Thus $N_c$ realises in $\mathrm{gr}(F)$ the
full Lie ideal $\bigoplus_{w\ge c}L_w$ generated by the degree-$c$ part; this is
the structural reason the quotient $Q$ is residually nilpotent. The passage from
this graded identification to residual nilpotence of $Q$ is the content of the
collection theory of free groups: in the notation of Magnus--Karrass--Solitar,
\emph{Combinatorial Group Theory}, \S5.3--5.7 and Theorem~5.13, and of Hall,
\emph{The Theory of Groups}, \S11.1--11.2, the collected normal form realises
$F/\gamma_m(F)$ as the free nilpotent group of class $m-1$ on the weight-$(<m)$
basic commutators, and a normal subgroup $N\trianglelefteq F$ whose graded image
in each degree $w$ is a \emph{direct summand} of $L_w$ (as is the case for
$N=N_c$, where the image is $L_w$ for $w\ge c$ and $0$ for $w<c$) is closed in the
lower central filtration, i.e.\ $N=\bigcap_w N\gamma_w(F)$. Applying this with
$N=N_c$ yields $\bigcap_w N_c\gamma_w(F)=N_c$, which is exactly
$\bigcap_k\gamma_k(Q)=\{1\}$.

\begin{lemma}\label{lem:gradedNc}
Let $\pi_w\colon\gamma_w(F)\to L_w=\gamma_w(F)/\gamma_{w+1}(F)$ be the projection.
For every $w\ge c$, $\pi_w(N_c\cap\gamma_w(F))=L_w$; for every $w<c$,
$N_c\subseteq\gamma_c(F)\subseteq\gamma_{w+1}(F)\cdot\gamma_c(F)$ contributes $0$
to $L_w$. Equivalently, the graded image of $N_c$ in $\mathrm{gr}(F)$ is
$\bigoplus_{w\ge c}L_w$.
\end{lemma}

\begin{proof}
For $w\ge c$: by Proposition~\ref{prop:moddescent}, inclusion \eqref{eq:E_w}, we
have $\gamma_w(F)\subseteq N_c\gamma_{w+1}(F)$. Given $g\in\gamma_w(F)$ write
$g=n t$ with $n\in N_c$, $t\in\gamma_{w+1}(F)$; then $n=g t^{-1}\in
\gamma_w(F)\cap N_c$ and $\pi_w(n)=\pi_w(g)$ (as $\pi_w(t)=0$). Hence
$\pi_w(N_c\cap\gamma_w(F))=L_w$. For $w<c$: $N_c\subseteq\gamma_c(F)\subseteq
\gamma_{w+1}(F)$, so $N_c\cap\gamma_w(F)=N_c\subseteq\gamma_{w+1}(F)$ maps to $0$
in $L_w$. This identifies the graded image of $N_c$ as claimed.
\end{proof}

\begin{remark}[On the status of the closedness statement]\label{rem:status}
Lemma~\ref{lem:gradedNc} identifies the graded image of $N_c$ exactly. The one
non-elementary step, used to pass from this graded identification to
$\bigcap_w N_c\gamma_w(F)=N_c$ (equivalently Proposition~\ref{prop:RN}), is the
closedness of $N_c$ in the lower central filtration. By Remark~\ref{rem:gap} this
closedness is logically equivalent, given part (II), to Theorem~\ref{thm:main}; it
is therefore not derivable from the graded data alone, and we invoke it as the
classical input recorded above (Magnus collection / Hall basis theory: a normal
subgroup whose graded image is, in each degree, a direct summand of $L_w$ is
closed in the lower central filtration of the free group). With that input,
Lemma~\ref{lem:gradedNc} gives $\bigcap_w N_c\gamma_w(F)=N_c$, hence
Proposition~\ref{prop:RN}, and Proposition~\ref{prop:final} then completes the
proof of Theorem~\ref{thm:main}.
\end{remark}

\bigskip
\noindent\textbf{Conclusion.}
Combining Lemma~\ref{lem:easy} (the inclusion $N_c\subseteq\gamma_c(F)$) with
Proposition~\ref{prop:final} (which yields $\gamma_c(F)\subseteq N_c$ from
Propositions~\ref{prop:moddescent} and~\ref{prop:RN}) gives
$\gamma_c(F)=N_c$ for every $c\ge1$. This proves Theorem~\ref{thm:main}: the
answer to the problem is \emph{yes}.

\end{document}
